% ---
% school: "college"
% level: "8eme"
% section: null
% chapter: "العمليات والحساب على الأعداد الكسرية"
% source: "Axiomaths"
% title: "تركيب: ميزانية أمين"
% country: "Tunisie"
% year: 2026
% professor: "Axiomaths"
% difficulty: "Moyen"
% points: null
% tags: ["fractions", "probleme", "budget", "synthese"]
% ---
% Exercice 7/15 de la série 8ème « العمليات والحساب على الأعداد الكسرية » (statement + \axsol + solution). Compiler avec XeLaTeX + axiomath.sty (+ axiomaths-logo.png).
\documentclass[12pt,a4paper]{article}
\usepackage[level=8b,color=15803D]{axiomath}
\axsetup{title={العمليات والحساب على الأعداد الكسرية}, corrige}

\begin{document}

\begin{corr}[title={تركيب: ميزانية أمين}, stars=3, source={Axiomaths}]
ادّخر أمين مبلغا من المال. أنفق $\dfrac{1}{3}$ هذا المبلغ في شراء كتب و $\dfrac{2}{5}$ منه في شراء لباس، وبقيت له $12$ دينارا.
\begin{enumerate}
  \item ما هو الكسر الذي يمثّله ما أنفقه أمين من المبلغ ؟
  \item ما هو الكسر الذي يمثّله ما بقي له ؟
  \item استنتج المبلغ الذي ادّخره أمين.
  \item احسب ثمن الكتب وثمن اللباس.
\end{enumerate}
\axsol
\begin{prop}[كسر من كمية]
$\dfrac{a}{b}$ من كمية $x$ هو الجذاء \ib{$\dfrac{a}{b}\times x$}
\end{prop}
\begin{enumerate}[itemsep=14pt]
  \item الكسر المنفق هو مجموع الكسرين :
  \[\dfrac{1}{3}+\dfrac{2}{5}=\dfrac{5}{15}+\dfrac{6}{15}=\dfrac{11}{15}\]
  أنفق أمين $\dfrac{11}{15}$ من المبلغ.
  \item المبلغ كله يمثل $1=\dfrac{15}{15}$ :
  \[\dfrac{15}{15}-\dfrac{11}{15}=\dfrac{4}{15}\]
  بقي لأمين $\dfrac{4}{15}$ من المبلغ.
  \item إذن $\dfrac{4}{15}$ من المبلغ يساوي $12$ دينارا.

  \textbf{طريقة التفكير :} $\dfrac{4}{15}$ تساوي $4$ أجزاء من $15$ جزءا متساويا، فهذه الأجزاء الأربعة تساوي $12$ دينارا.
  الجزء الواحد يساوي $12\div4=3$ دنانير.

  المبلغ كله يساوي $15$ جزءا، أي $15\times3=45$ دينارا.

  وبالقسمة على كسر نجد نفس النتيجة : \ib{$12\div\dfrac{4}{15}=12\times\dfrac{15}{4}=45$}
  \item ثمن الكتب :
  \[\dfrac{1}{3}\times45=15\]
  أي $15$ دينارا. ثمن اللباس :
  \[\dfrac{2}{5}\times45=18\]
  أي $18$ دينارا.
  \textbf{للتأكد :} $15+18+12=45$ $\checkmark$
\end{enumerate}
\textbf{الفكرة :} نحسب أولا الكسر الباقي ثم نعود إلى الكمية الكلية بالقسمة على هذا الكسر.
\end{corr}

\end{document}
